% slides/modules/decision_theory_metrics.tex
% Cap. 5 (extracto) — Decision Theory: lo esencial para métricas

% -------------------- Slide: Section title --------------------
\section{Decision Theory para métricas}

% -------------------- Slide 1 --------------------
\begin{frame}{Motivación: de creencias a decisiones}
\begin{itemize}
  \item La inferencia Bayesiana actualiza creencias: \(\;p(H\mid x)\).
  \item Pero al final necesitamos una \textbf{acción} \(a \in \mathcal{A}\) (ej.\ diagnosticar, aprobar crédito, detectar fraude).
  \item La \textbf{función de pérdida} \(\ell(h,a)\) cuantifica el costo de tomar \(a\) cuando el estado real es \(h\).
  \item \textbf{Idea clave:} elegir la acción que minimiza la pérdida esperada bajo el posterior.
\end{itemize}
\end{frame}

% -------------------- Slide 2 --------------------
\begin{frame}{5.1.1 Basics: riesgo (pérdida esperada) y regla óptima}
\textbf{Riesgo posterior (pérdida esperada):}
\[
\Risk(a\mid x) \;\triangleq\; \E_{p(h\mid x)}[\ell(h,a)]
\;=\;\sum_{h\in \mathcal{H}} \ell(h,a)\,p(h\mid x)
\]

\textbf{Política óptima (regla de decisión Bayesiana):}
\[
\pi^*(x) \;=\;\argmin_{a\in \mathcal{A}} \Risk(a\mid x)
\]

\textbf{Equivalente en utilidad:} si \(U(h,a)=-\ell(h,a)\),
\[
\pi^*(x) \;=\;\argmax_{a\in \mathcal{A}} \E[U(h,a)]
\]

\vspace{0.4em}
\begin{itemize}
  \item Conecta directamente con \textbf{métricas}: distintos costos \(\Rightarrow\) distintas reglas/umbrales.
\end{itemize}
\end{frame}

% -------------------- Slide 3 --------------------
\begin{frame}{Clasificación: 0-1 loss \(\Rightarrow\) MAP}
Supón clases \(Y=\{1,\dots,C\}\) y acciones \(A=Y\).

\textbf{Pérdida 0-1:}
\[
\ell_{01}(y^*,\hat y)=\I(y^*\neq \hat y)
\]

\textbf{Riesgo:}
\[
\Risk(\hat y\mid x)=P(\hat y\neq y^*\mid x)=1-P(y^*=\hat y\mid x)
\]

\textbf{Decisión óptima:}
\[
\pi(x)=\argmax_{y\in Y} p(y\mid x)
\]
\begin{itemize}
  \item Esta es la regla \textbf{MAP} (escoger la clase más probable).
  \item \textbf{Limitación:} asume que todos los errores cuestan igual.
\end{itemize}
\end{frame}

% -------------------- Slide 4 --------------------
\begin{frame}{Cost-sensitive classification: umbral por costos}
Binario \(y\in\{0,1\}\). Matriz de pérdidas:
\[
\begin{pmatrix}
\ell_{00} & \ell_{01}\\
\ell_{10} & \ell_{11}
\end{pmatrix}
\]

Sea \(p_1=p(y^*=1\mid x)\), \(p_0=1-p_1\).
Elegimos \(\hat y=0\) si:
\[
\ell_{00}p_0+\ell_{10}p_1 \;<\; \ell_{01}p_0+\ell_{11}p_1
\]

Caso común \(\ell_{00}=\ell_{11}=0\):
\[
 p_1 \;<\; \frac{\ell_{01}}{\ell_{01}+\ell_{10}}
\]

\begin{itemize}
  \item Si \(\ell_{10}=c\,\ell_{01}\) (FN cuesta \(c\) veces FP), entonces
  \[
  p_1 < \frac{1}{1+c}
  \]
  \item \textbf{Interpretación:} el \textbf{umbral} no es “mágico”; viene de costos/objetivo del negocio.
\end{itemize}
\end{frame}

% -------------------- Slide 5 --------------------
\begin{frame}{Opción de rechazo: “no sé” con costo controlado}
Acciones \(A=Y\cup\{0\}\), donde \(a=0\) es \textbf{rechazar}.
\[
\ell(y^*,a)=
\begin{cases}
0 & \text{si } a=y^* \in \{1,\dots,C\}\\
\lambda_r & \text{si } a=0\\
\lambda_e & \text{si } a\neq y^*,\ a\in \{1,\dots,C\}
\end{cases}
\]

\textbf{Regla óptima (Chow):}
\[
a^*=
\begin{cases}
 y^* & \text{si } p^*>\lambda^*\\
\text{rechazar} & \text{en otro caso}
\end{cases}
\quad\text{donde}\quad
 y^*=\argmax_y p(y\mid x),\; p^*=\max_y p(y\mid x),\;
\lambda^*=1-\frac{\lambda_r}{\lambda_e}
\]

\begin{itemize}
  \item Útil en dominios de alto riesgo (medicina, finanzas).
  \item Conecta con \textbf{calibración} y \textbf{confianza} del modelo.
\end{itemize}
\end{frame}

% -------------------- Slide 6 --------------------
\begin{frame}{De decisiones a métricas: matriz de confusión}
Para un umbral fijo, obtenemos: TP, FP, TN, FN.

\begin{center}
\begin{tabular}{c ccc}
\toprule
 & \multicolumn{2}{c}{\textbf{Predicción}}\\
\cmidrule(lr){2-3}
\textbf{Real} & 0 & 1\\
\midrule
0 & TN & FP\\
1 & FN & TP\\
\bottomrule
\end{tabular}
\end{center}

\textbf{Tasas (normalización por fila):}
\[
\text{TPR (Recall)}=\frac{TP}{TP+FN},\quad
\text{FPR}=\frac{FP}{FP+TN},\quad
\text{FNR}=1-\text{TPR}
\]

\textbf{Por columna:}
\[
\text{Precision}=\frac{TP}{TP+FP},\quad
\text{NPV}=\frac{TN}{TN+FN}
\]

\begin{itemize}
  \item Estas cantidades dependen del \textbf{umbral} \(\tau\).
\end{itemize}
\end{frame}

% -------------------- Slide 7 --------------------
\begin{frame}{ROC curve: TPR vs FPR al variar el umbral}
\begin{itemize}
  \item Un clasificador produce un \textbf{score} (idealmente \(p(y=1\mid x)\)).
  \item Para cada umbral \(\tau\): obtenemos un punto \((\text{FPR}(\tau),\text{TPR}(\tau))\).
  \item La curva ROC es el trazo de \(\text{TPR}\) contra \(\text{FPR}\) al variar \(\tau\).
\end{itemize}

\vspace{0.4em}
\textbf{Resumen escalar:}
\begin{itemize}
  \item \textbf{AUC-ROC}: área bajo la ROC (más alto es mejor).
  \item \textbf{EER}: punto donde \(\text{FPR}=\text{FNR}\) (más bajo es mejor).
\end{itemize}

\vspace{0.4em}
\textbf{Nota sobre desbalance:}
\begin{itemize}
  \item ROC es relativamente insensible al desbalance, pero puede ocultar muchos FP absolutos.
\end{itemize}
\end{frame}

% -------------------- Slide 8 --------------------
\begin{frame}{Precision-Recall curve: mejor para clases raras}
\textbf{Definiciones (para umbral \(\tau\)):}
\[
P(\tau)=\frac{TP_\tau}{TP_\tau+FP_\tau},\qquad
R(\tau)=\frac{TP_\tau}{TP_\tau+FN_\tau}
\]

\begin{itemize}
  \item PR curve: graficar \textbf{Precision} vs \textbf{Recall} al variar \(\tau\).
  \item En problemas con \textbf{clase positiva rara}, PR suele ser más informativa que ROC.
\end{itemize}

\vspace{0.4em}
\textbf{Resumen escalar:}
\begin{itemize}
  \item \textbf{Average Precision (AP)}: promedio de precisión interpolada (relacionado con área bajo PR interpolada).
  \item \textbf{mAP}: promedio de AP sobre múltiples curvas/tareas/clases.
\end{itemize}
\end{frame}

% -------------------- Slide 9 --------------------
\begin{frame}{F-scores: combinar Precision y Recall}
Para un umbral fijo (un punto en la PR curve), define \(P\) y \(R\).

\textbf{F\(_\beta\):}
\[
F_\beta \;=\; (1+\beta^2)\frac{P\,R}{\beta^2 P + R}
\]

\textbf{Caso \(\beta=1\) (armoniza ambos):}
\[
F_1=\frac{2PR}{P+R}
\]

\begin{itemize}
  \item Media armónica penaliza estrategias que “ganan” solo en \(P\) o solo en \(R\).
  \item \(\beta>1\): prioriza \textbf{recall}; \(\beta<1\): prioriza \textbf{precision}.
\end{itemize}
\end{frame}

% -------------------- Slide 10 --------------------
\begin{frame}{Efecto del desbalance: ROC vs PR (intuición)}
Sea \(\pi=\frac{P}{P+N}\) (prevalencia), \(r=\frac{P}{N}=\frac{\pi}{1-\pi}\).

\textbf{ROC:} \(\text{TPR}\) y \(\text{FPR}\) son tasas internas a positivos/negativos \(\Rightarrow\) menos sensibles a \(\pi\).

\textbf{PR:} la precision depende de \(\pi\):
\[
\text{Prec}=\frac{TP}{TP+FP}
=\frac{\text{TPR}}{\text{TPR}+\frac{1}{r}\text{FPR}}
\]
\begin{itemize}
  \item Si \(\pi\) baja (positivos raros), la precision tiende a bajar para los mismos \(\text{TPR},\text{FPR}\).
  \item Por eso \textbf{PR-AUC / AP} es muy usada en riesgo/fraude/default.
\end{itemize}
\end{frame}

% -------------------- Slide 11 --------------------
\begin{frame}{Takeaways para Semana 1}
\begin{itemize}
  \item Métricas \(\approx\) \textbf{pérdidas implícitas}: elige según el costo de errores (FP vs FN).
  \item La decisión óptima depende de \(p(y\mid x)\) \(\Rightarrow\) importa \textbf{calibración} y \textbf{umbral}.
  \item ROC: útil para comparar rankings; PR: preferible con \textbf{clase positiva rara}.
\end{itemize}

\vspace{0.6em}
\centering
\textit{Regla práctica: si el evento positivo es raro y caro, mira PR-AUC/AP y costos.}
\end{frame}
